\devoirDS{3}{Chapitres 7 à 10}{Une réserve de glace/6,Détartrage contrôlé/7,Trois transformations/7}
\begin{blocLivre}{ds-3-1}\begin{exoBac}{6}{Une réserve de glace}
Une enceinte contient 150 g de glace à 0 °C. $L_f=334$ kJ/kg. Le milieu lui apporte utilement 50 J par seconde, supposés constants.
\begin{questions}
\item Écrire la fusion de l’eau et qualifier son caractère thermique.
\item Calculer l’énergie nécessaire à la fusion complète.
\item Estimer la durée et expliquer ce qui change ensuite.
\end{questions}
\end{exoBac}\end{blocLivre}
\begin{blocLivre}{ds-3-2}\begin{exoBac}{7}{Détartrage contrôlé}
On fait réagir 0,040 mol de \ce{CaCO3} et 0,060 mol de \ce{H+}. L’équation est $\ce{CaCO3 + 2 H+ -> Ca^{2+} + CO2 + H2O}$. Des ions chlorure restent inchangés.
\begin{questions}
\item Nommer l’espèce spectatrice.
\item Déterminer le réactif limitant.
\item Donner les quantités de calcaire restant et de gaz produit.
\item Proposer un test du gaz et une précaution concernant le récipient.
\end{questions}
\end{exoBac}\end{blocLivre}
\begin{blocLivre}{ds-3-3}\begin{exoBac}{7}{Trois transformations}
Un laboratoire compare un produit synthétisé S à une référence N et un réactif R sur la plaque fournie. Un autre document donne $\ce{^{2}_{1}H + ^{3}_{1}H -> ^{4}_{2}He + ^{1}_{0}n}$.
\schema{ccm-synthese}
\begin{questions}
\item Quelle différence fondamentale sépare synthèse chimique et réaction nucléaire ?
\item Conclure sur l’identité probable et la pureté de S à partir de la CCM.
\item Expliquer le rôle du réfrigérant d’un reflux utilisé pour la synthèse.
\item Vérifier les bilans nucléaires et nommer la réaction donnée.
\end{questions}
\end{exoBac}\end{blocLivre}

